\subsection{How the design is checked}
\label{sec:ha:enabled:model_checking}

The scenarios in \texttt{scripts/ha\_test.py} run the venue and inject failures.
Each scenario tries one sequence of events. The failures that matter most in this
section depend on the order in which events happen: a message arriving just
after a restart, or two timeouts firing close together. A scenario that passes
says nothing about the orders it did not try.

\textbf{The design is therefore also checked by model checking.} The design is
written down in TLA+, a language for describing a system precisely enough for a
program to check it. The description states what each component holds, which
steps can change it, and which properties must always hold. The model checker,
TLC, then visits every state that a small version of the system can reach, in
every order of events, and reports any sequence of steps that breaks a property.
That sequence is called a counterexample.

The check has two limits. It covers a small version of the system: two instances,
one arbiter or two, and one or two failures at a time. It checks the design as
written in TLA+, so each property it confirms holds in the code only as far as
the TLA+ description matches the code.

\subsubsection{What is checked}
\label{sec:ha:enabled:model_checking:scope}

Two TLA+ descriptions model the design the code implements:

\begin{itemize}
  \item the two sequencer instances and the arbiter they report to; and
  \item the two arbiters and the witness, which decide which arbiter is active.
\end{itemize}

Both include crashes and restarts of every party, the failure of any link, and
messages that arrive after the state they describe has changed.

Model checking found eight defects in that design. The code now addresses six of
them in full and part of one more. The properties that now hold within the limits
checked are these:

\begin{itemize}
  \item no two leaders ever lead in the same epoch, so a receiver can always tell
        them apart;
  \item two leaders that can hear each other always resolve to one; and
  \item two active arbiters that can hear each other always resolve to one.
\end{itemize}

\textbf{Two instances can still both lead while a partition lasts.} An instance
that can reach no arbiter promotes itself, as \reqref{R-0095} permits. When the
real cause is that the instance has been cut off, the pair has two leaders until
the partition heals. The arbiter pool can likewise have two active arbiters during
a partition, because it does not yet decide by majority.

A replacement design removes promotion without an arbiter. In the replacement, an
instance leads only while a majority of three voters has granted it a lease that
has not run out. It is a proposal, described in
\texttt{docs/availability/majority\_leases.md}, and nothing in the code implements
it. A third TLA+ description specifies it. Model checking that description found
no state in which two instances act as leader at once. It also showed that each
rule of the replacement is needed: with any one rule removed, TLC finds two
instances acting as leader at once within a few steps.

The descriptions, the counterexamples and the full results are in
\texttt{docs/availability/tla/}, with the results in \texttt{findings.md}.

\subsubsection{What the venue's build must do}
\label{sec:ha:enabled:model_checking:behaviour}

\begin{requirement}{R-0145}{The build reproduces every recorded counterexample}
  The build shall rerun each counterexample recorded against a TLA+ description of
  the high availability design, and shall fail if any of them no longer breaks the
  property it is recorded against.
  \rationale{A counterexample is the evidence that a rule of the design is needed.
  A later change to the TLA+ description can stop a counterexample failing, and
  the description then no longer shows what the recorded results say it shows.
  Nothing else in the build would notice. Rerunning every counterexample takes a
  few seconds.}
  \covers{tla:lease-1-holder-counts-from-arrival, tla:lease-2-candidate-ignores-own-grant,
  tla:lease-3-restart-does-not-wait, tla:lease-4-degraded-promotion,
  tla:lease-5-epoch-regresses, tla:lease-6-no-yield-livelock, tla:lease-7-no-peer-vote}
\end{requirement}

The counterexamples reproduced today are those of the replacement design.
\texttt{scripts/tla\_trace\_pages.py} reruns them at install time, and writes a web
page that steps through each one.
